\documentclass[10pt]{article}
\usepackage[margin=0.8in]{geometry}
\usepackage[]{xcolor}
\usepackage{amsmath}
\parindent=0pt

\title{CSE 4300: Operating Systems Notes}
\author{gordonbchen}
\date{}

\begin{document}
\maketitle


\section{Computers}

\subsection{Boot}
\begin{enumerate}
    \item Power on.
    \item CPU execute BIOS (firmware on Motherboard).
    \item BIOS loads bootloader into RAM, starts executing bootloader.
    \item Bootloader loads OS into RAM and starts executing OS.
\end{enumerate}

\subsection{Threads vs Processes}
Process: resource container (own virtual address space), has multiple threads.

Thread: unit of execution, has PC, registers, stack.

Core = CPU. $n$-core CPU can execute $n$ threads in parallel.

Context switching simulates more cores: suspends running proc on interrupt, OS schedules next proc.\\

Syscall
\begin{itemize}
    \item User program wants OS service.
    \item Switch from user to kernel mode.
    \item OS runs, gives back control after syscall.
\end{itemize}


\section{Processes}
Process table
\begin{itemize}
    \item PID: process ID
    \item State: running, blocked, ready
    \item Dumped state: PC, registers
    \item Memory: files, memory addr
\end{itemize}

Process memory
\begin{itemize}
    \item Stack: function call (stack frame), grows down
    \item Heap: dynamically allocated, malloc, grows up
    \item Data: global and static vars
    \item Code
\end{itemize}

Context switch
\begin{enumerate}
    \item Timer interrupt / syscall
    \item Hardware pushes current PC onto stack, jumps to addr in interrupt vector table
    \item Asm dumps PC and registers into Process Table and creates new stack for interrupt
    \item Interrupt runs
    \item Scheduler runs next process
\end{enumerate}

$\text{CPU util} = 1 - (\text{IO wait \%})^{\text{\# procs}}$


\section{Threads}
threads vs processes
\begin{itemize}
    \item pthread\_create vs fork
    \item threads in the same process share memory, every process has different virtual address space
\end{itemize}

User-level threads (JVM)
\begin{itemize}
    \item process table in kernel space
    \item each process is a runtime with a thread table and threads in user-space
    \item problem: user-space thread wants to read file, OS blocks the entire process even if
        other threads can run.
\end{itemize}

Kernel threads: process and thread tables manages by kernel

CPU quantum: default time process runs for, can be interrupted


\section{Race Conditions}
Peterson's Solution and XCHG: busy waiting, always tries to relock\\
Prone to priority inversion problem: L never scheduled to unlock b/c H is busy waiting

\begin{verbatim}
mutex_lock:
    TSL register, mutex  // copy mutex to register, set mutex to 1. atomic, locks mem bus.
    CMP register, #0     // compare register to 0 (did we get the mutex)
    JMP ok               // if mutex 0, then we have it, return to user code
    CALL thread_yield    // mutex is 1, we did not get the mutex, tell OS we are waiting
    JMP mutex_lock       // we got woken up, try to relock the mutex
ok:
    RET

mutex_unlock:
    MOVE mutex, #0       // set the mutex to 0
    RET
\end{verbatim}

Exponential backoff when multiple threads try to access the same resource.

Reader / writer problem
\begin{itemize}
    \item Writer need exclusive access to write to database, just tries to lock.
    \item First reader locks db, last reader unlocks db, track \# readers
\end{itemize}

Other primitives: condition vars (wait until condition and signal condition) and barriers (synchronize threads)


\section{Scheduling}
Batch (slurm)
\begin{itemize}
    \item First come first serve: FIFO queue, blocked proc gets moved to the end of the queue
    \item Shortest first: optimal turnaround time if all jobs arrive at the same time
    \item Shortest remaining time next: pre-emptive shortest first
\end{itemize}

Interactive (MS Word)
\begin{itemize}
    \item Round-robin
    \item Priority: each proc assigned priority, reduce priority to not starve low priority procs\\
        Higher priority to IO bound: priority = 1 / frac of quantum used\\
        Round-robin in priority classes
    \item Multiple queues: exponetially increasing quantum allocation until proc finishes (less swaps)
    \item Shortest process next: aging $T_n = \alpha T_{n-1} + (1-\alpha) T_{n}$
    \item Guarenteed scheduling: each user gets $1/n$ CPU share, track proc time
    \item Lottery scheduling: randomized, users buy lottery tickets to run procs
\end{itemize}

Real-time\\
Earliest deadline first: $\sum \dfrac{\text{cpu time}}{\text{period time}} < 1$

\section{Memory Management}
\subsection{Relocation problem}
Compilers assume programs are loaded at mem addr 0, but they can be loaded anywhere in RAM.

Static relocation: add program start addr to all addresses (slow, must be recomputed each load)

Base and limit registers: hardware adds base register to addr, checks under limit register

\subsection{Swapping}
Move procs back and forth from HDD.

Memory holes: compact into contiguous, leave space for procs to dynamically allocate

Bitmap: each RAM chunk is bit in bitmap

Linked list: proc / hole start, length, pointer to next

First fit is best.

\subsection{Virtual Memory}
MMU (Memory Management Unit): translates virtual mem addr to physical meme addr\\

Page table
\begin{itemize}
    \item virtual memory addr = page table idx + offset
    \item \# bits in offset determined by page size
    \item physical memory addr = page\_table[page\_table\_idx] + offset
    \item page fault: present bit = 0 in page table, page not loaded in RAM, must fetch from HDD.
    \item dirty (modified) bit: must write-back to RAM if removed from page table
    \item referenced bit: set if recently accessed, periodically cleared
    \item protections: RWX permissions
\end{itemize}

Memory access
\begin{itemize}
    \item virtual addr = page table idx + offset
    \item MMU checks page\_table[page\_idx] present bit
    \item if present bit = 0, page fault, load page from HDD into RAM, writing back if existing page dirty
    \item physical addr = physical page + offset
\end{itemize}

TLB: Translation Lookaside Buffer
\begin{itemize}
    \item cache of page table in MMU (entire page table cannot fit)
    \item TLB miss: page not in TLB, go look in RAM page table\\
        Soft miss: page loaded, found in RAM table\\
        Hard miss: page not loaded, not in RAM table, go fetch from HDD
    \item Context switch: flush TLB b/c each proc has own virtual address space
\end{itemize}

Multilevel page tables
\begin{itemize}
    \item virtual addr = level 1 page idx + level 2 page idx + offset
    \item only create level 2 page tables you need
\end{itemize}

Inverted page tables: hash virtual page $\to$ linked list of (virtual page, process id, physical page)

\subsubsection{Page replacement algorithms}
FIFO: throw out oldest page

Second chance: FIFO but if referenced bit set, put to end of queue\\


Clock page replacement algorithm: 2nd chance w/ circular buffer, go around until no R-bit

NRU (Not Recently Used): discard order (R, M) = (0, 0), (0, 1), (1, 0), (1, 1)\\


LRU (Least Recently Used): increment counter on reference, take lowest counter

NFU (Not Frequently Used): add R-bit to per-page counter, replace page with lowest counter, doesn't forget (bad)

Aging alg: add R-bit to left of per-page counter, replace lowest, 8-bit memory, can't know in a tick who was accessed first\\


Working set
\begin{itemize}
    \item in page table: track time of last use for page and referenced bit
    \item if referenced: set time of last use to current time
    \item else: throw away page with smallest time
\end{itemize}

Working set clock
\begin{itemize}
    \item old page and unmodified: claim
    \item old page and modified: schedule write
    \item after going all the way around:\\
        if write scheduled: keep scanning, will be clean page to claim\\
        else: all pages in working set, claim any
\end{itemize}


\section{Exam 1 Review Problems}

\end{document}
